% 广西自然科学基金面上项目报告正文项目级配置。
% 本项目独立于 bensz-nsfc，避免影响 NSFC 与其它省级基金模板。

\usepackage{xcolor}
\usepackage{graphicx}
\usepackage{amsmath}
\usepackage{amssymb}
\usepackage{bm}
\usepackage{geometry}
\usepackage{fontspec}
\usepackage{xeCJK}
\usepackage{ragged2e}
\usepackage{indentfirst}
\usepackage{enumitem}
\usepackage{caption}
\usepackage{booktabs}
\usepackage{tabularx}
\usepackage{longtable}
\usepackage{array}
\usepackage{hyperref}

% Word 基线：A4，上下 1440 twips（2.54 cm），左右 1800 twips（3.175 cm）。
\geometry{a4paper,left=3.175cm,right=3.175cm,top=2.54cm,bottom=2.54cm}

\hypersetup{
  colorlinks=true,
  urlcolor=black,
  linkcolor=black,
  anchorcolor=black,
  citecolor=black,
  CJKbookmarks=true
}

\newcommand{\GXNSFFallbackFontPath}{}

\newcommand{\GXNSFSetupLocalFontFiles}[1]{%
  \setmainfont[Path={#1},Extension=.ttf,AutoFakeBold=2.5,AutoFakeSlant=0.2]{TimesNewRoman}%
  \renewcommand{\GXNSFFallbackFontPath}{#1}%
}

\newcommand{\GXNSFSetupInstalledFonts}{%
  \usepackage{bensz-fonts}%
  \BenszFontsSetMainTimesFallback%
  \renewcommand{\GXNSFFallbackFontPath}{\BenszFontsDir}%
}

\newcommand{\GXNSFSetupFonts}{%
  \IfFileExists{styles/fonts/AdobeFangsongStd-Regular.otf}{%
    \GXNSFSetupLocalFontFiles{styles/fonts/}%
  }{%
    \IfFileExists{bensz-fonts.sty}{%
      \GXNSFSetupInstalledFonts
    }{%
      \PackageError{GXNSF}{Missing required package bensz-fonts}{%
        Install bensz-fonts with the repository-level scripts/install.py before compiling the standard package.\MessageBreak
        The Overleaf release already embeds the required fonts.%
      }%
    }%
  }%
  \GXNSFPreferOriginalFonts
  \xeCJKsetup{PunctStyle=plain,CJKecglue={}}%
}

% 原 DOCX 指定“方正仿宋_GBK / 方正楷体_GBK”。优先使用系统已安装的
% 对应方正字体；若不存在，则保留上面的仓库内置字体，以保证跨平台可编译。
\newcommand{\GXNSFSetSystemFangsong}[1]{%
  \setCJKmainfont[AutoFakeBold=2.5]{#1}%
  \setCJKfamilyfont{gxnsffangsong}[AutoFakeBold=2.5]{#1}%
}

\newcommand{\GXNSFSetSystemKai}[1]{%
  \setCJKfamilyfont{gxnsfkai}[AutoFakeBold=2.5]{#1}%
}

\newcommand{\GXNSFSetFallbackFangsong}{%
  \setCJKmainfont[Path={\GXNSFFallbackFontPath},AutoFakeBold=2.5]{AdobeFangsongStd-Regular.otf}%
  \setCJKfamilyfont{gxnsffangsong}[Path={\GXNSFFallbackFontPath},AutoFakeBold=2.5]{AdobeFangsongStd-Regular.otf}%
}

\newcommand{\GXNSFSetFallbackKai}{%
  \setCJKfamilyfont{gxnsfkai}[Path={\GXNSFFallbackFontPath},AutoFakeBold=2.5]{Kaiti.ttf}%
}

\newcommand{\GXNSFPreferOriginalFonts}{%
  \IfFontExistsTF{方正仿宋_GBK}{%
    \GXNSFSetSystemFangsong{方正仿宋_GBK}%
  }{%
    \IfFontExistsTF{方正仿宋简体}{%
      \GXNSFSetSystemFangsong{方正仿宋简体}%
    }{\GXNSFSetFallbackFangsong}%
  }%
  \IfFontExistsTF{方正楷体_GBK}{%
    \GXNSFSetSystemKai{方正楷体_GBK}%
  }{%
    \IfFontExistsTF{方正楷体简体}{%
      \GXNSFSetSystemKai{方正楷体简体}%
    }{\GXNSFSetFallbackKai}%
  }%
}

\GXNSFSetupFonts

\newcommand{\gxnsffangsong}{\CJKfamily{gxnsffangsong}}
\newcommand{\gxnsfkai}{\CJKfamily{gxnsfkai}}
\newcommand{\GXNSFFontSize}{\fontsize{16pt}{28.3pt}\selectfont}

\pagestyle{empty}
\renewcommand{\baselinestretch}{1.0}
\setlength{\parindent}{32pt}
\setlength{\parskip}{0pt}
\XeTeXlinebreaklocale "zh"
\XeTeXlinebreakskip=0pt plus 1pt
\AtBeginDocument{\GXNSFFontSize}

\captionsetup{font={small},labelsep=period,labelfont=bf}
\setlist{nosep,leftmargin=2em}

\newcommand{\GXNSFTemplateParagraph}[2][]{%
  \par
  \begingroup
    \justifying
    \setlength{\parindent}{32pt}%
    \setlength{\parskip}{0pt}%
    \GXNSFFontSize
    #1#2\par
  \endgroup
}

\newcommand{\GXNSFReportTitle}{%
  % 锚定首页 top skip，使标题可见上边界与 Word 导出 PDF 的约 2.85 cm 对齐。
  \hrule height 0pt width 0pt\relax
  \vskip -2pt\relax
  \GXNSFTemplateParagraph[\centering\setlength{\parindent}{0pt}\gxnsffangsong\bfseries]{报告正文}%
}

\newcommand{\GXNSFPart}[1]{%
  \GXNSFTemplateParagraph[\gxnsfkai]{#1}%
}

\newcommand{\GXNSFItem}[2]{%
  \GXNSFTemplateParagraph[\gxnsffangsong]{{\bfseries #1}#2}%
}

\newcommand{\GXNSFItemMixed}[3]{%
  \GXNSFTemplateParagraph[\gxnsffangsong]{{\bfseries #1}#2{\bfseries #3}}%
}

\newcommand{\GXNSFBodyText}{%
  \setlength{\parindent}{32pt}%
  \setlength{\parskip}{0pt}%
  \GXNSFFontSize
  \gxnsffangsong
}

\pdfstringdefDisableCommands{%
  \def\gxnsffangsong{}%
  \def\gxnsfkai{}%
}
